\section{OpenHands 平台与 Agent 架构}
\subsection{OpenHands 架构组件}
OpenHands 提供了一个容器化、可编排的 Agent 开发平台：Agent 能在受控的 sandbox 环境里同时操作终端、编辑器和浏览器，后端可以接入多个 LLM（如 code-davinci 及开源模型），并通过 webhook、测试解释器等工具链完成输入→推理→行动的闭环。

各组件通过事件总线和消息队列通信，允许在运行时替换模型或工具。平台还为每一次动作定义了元数据（触发方式、上下文、预期输出），方便做回滚与调试。

\subsection{多模态接口与审计}
\begin{figure}[H]
\centering
\includegraphics[width=\textwidth]{frame-01-openhands.jpg}
\caption{OpenHands 画面展示了 Agent 级联的环境感知、LLM 推理与工具执行流。 \protect\footnotemark}
\end{figure}
\footnotetext{视频画面时间区间：00:04:45--00:04:50。}

平台强调“可观察性”：每个 Agent 的行动（终端命令、浏览器点击、文件修改）都会被记录，测试结果与日志会流回一个中央审计台。OpenHands 同时提供接口供人工阻断、利用调试视图或插入辅助脚本，这种 plug-and-play 的能力让 Agent 不再是黑盒。

\begin{knowledgebox}{OpenHands 核心能力}
OpenHands 包含多模态接入器、LLM 决策引擎、工具执行器（Docker shell/测试/浏览器）以及可视化审计面板。各部分通过事件总线通信，可在运行时替换模型或工具。
\end{knowledgebox}

\begin{warningbox}{可观察性与安全边界}
Agent 记录每一个命令，但也必须限制高权限操作，确保沙箱可还原，否则开发日志与环境状态会泄露敏感信息、难以复现。
\end{warningbox}

\subsection{本章小结}
\begin{itemize}
    \item OpenHands 把 Agent 的感知、推理、行动与审计串成可插拔的流水线。
    \item 透明日志与人工干预接口是让工程团队愿意交付 Agent 的关键保障。
\end{itemize}

